\section{Neighbours and nomenclature}\label{sec:related}

\subsection{Trademarks as innovation data}

A trademark filing is a dated, costly, public disclosure made by the
millions of firms that never patent, and a literature has grown around
reading it as an innovation indicator \citep{Castaldi2020, Willeke2026}.
Trademark counts and breadth predict venture valuations
\citep{Block2014}; filing timing tracks innovation processes
\citep{Seip2018}; and trademarks mark the product-market side of
innovation that patents miss \citep{Flikkema2019, Flikkema2025}.
Census-linked work reads registration as a correlate of firm growth and
survival \citep{Dinlersoz2018, Dinlersoz2023}. Nearly all of it counts
filings, classes, or renewals; the descriptions' text enters, if at all,
as breadth or similarity. This paper reads the text itself, and reads it
against time.

\subsection{Text measures of novelty}

Scoring a document by its divergence from a reference corpus before and
after it comes from cultural analytics: \citet{Murdock2017} on Darwin's
reading notebooks, \citet{Barron2018} on the French Revolution's
parliamentary debates. In economics and finance, text-based measures
locate firms in product space \citep{HobergPhillips2016}, date
technological innovation in patents over the long run \citep{Kelly2021},
and track the vintage of patent language \citep{Kalyani2024}. The
construction here differs in its object --- the composition of commercial
offerings, not ideas or claims --- and in its two-sided reference: every
filing is scored against its own industry's language both before and
after its date, which is what lets \emph{unusual} be separated from
\emph{early}. The separation matters because the two components predict
opposite things.

\subsection{Nomenclature}\label{sec:nomenclature}

The innovation literature is not consistent about what quantities like
these are called, and several of its established terms assert things this
design cannot show, so the paper's names are chosen against the nearest
occupants rather than borrowed from them. Table~\ref{tab:nomenclature}
maps the five quantities; the entries below say, for each declined term,
what adopting it would claim.

\begin{table}[!htbp]\centering\small
\caption{The five quantities, their nearest occupants in the literature,
and the terms used here. $P$ is a filing's theme distribution;
the references $Q^{-}$ and $Q^{+}$ aggregate the five years of
same-class filings before and after it (Section~\ref{sec:measure}).
\emph{Lagging} and \emph{leading} describe the two signs of $L$. The
level names appear only where the two-sided construction itself is at
issue; the analysis runs on $A$ and $L$.}
\label{tab:nomenclature}
\begin{tabular}{p{3.4cm}p{5.6cm}p{4.4cm}}
\toprule
Quantity & Nearest in the literature & Term used here \\
\midrule
$K^{-} = \mathrm{KL}(P \| Q^{-})$ & novelty \citep{Murdock2017, Barron2018}; originality \citep{HallJaffe2001} & past-facing surprise \\
$K^{+} = \mathrm{KL}(P \| Q^{+})$ & transience \citep{Murdock2017}; generality \citep{HallJaffe2001}, which measures realized influence, not position & future-facing surprise \\
$A = \tfrac{1}{2}(K^{-} + K^{+})$ & atypicality \citep{Uzzi2013}, which scores internal pairings, not distance from the class & \textbf{atypicality} \\
$L = K^{-} - K^{+}$ & resonance \citep{Barron2018}, which asserts a mechanism & \textbf{lead} (leading / lagging) \\
$|L|$ & no established occupant & lead magnitude \\
\bottomrule
\end{tabular}
\end{table}

\emph{Novelty} and \emph{transience} are the two KL levels themselves;
novelty is declined as a name for anything here because it names only the
backward-looking half of a quantity that faces both ways.
\emph{Resonance} is arithmetically what this paper calls lead, and is
avoided because it names a mechanism --- the field resonating with an
author's contribution --- that this design cannot observe: a filing can
sit where its class was heading because it was imitated, because both
responded to the same outside event, or by coincidence.

\emph{Radicalness} is the closest neighbour among the innovation labels,
since KL is a departure measure rather than a newness count. It is
declined because it claims existing competences are devalued, which
nothing here licenses --- and because the label itself is unsettled: a
bibliometric review of 4{,}407 articles finds \emph{radical},
\emph{breakthrough}, \emph{discontinuous} and \emph{disruptive} used with
definitions inconsistent within labels and not distinctive across them
\citep{Verhoeven2025}.

\emph{Atypicality in the bibliometric sense} \citep{Uzzi2013}, applied to
patents as unconventionality by \citet{BerkesGaetani2021}, scores the
rarity of the pairwise combinations inside a document against the
population's co-occurrence rates. The atypicality here is a different
object: the distance of one document's composition from its class's, not
the rarity of its internal pairings; the two are close to independent on
this corpus (Section~\ref{sec:measure}). The word is kept because both
mean unlike the population; which population, and along what, differ.

\emph{Originality} and \emph{generality} \citep{HallJaffe2001} are the
patent literature's standard backward- and forward-facing pair, and
generality is the nearest established forward-facing quantity to $K^+$.
Both measure the dispersion of citations across technology classes,
and generality requires the forward citations to have arrived, so it
observes realized influence; $K^+$ observes the class's movement whether
or not the filing influenced it.

\emph{Distinctiveness} has two prior occupants: trademark law's
registrability doctrine, and strategic management's optimal
distinctiveness \citep{Zhao2017}, which predicts an inverted-U that
appears here at registration but not at listing, so the term would imply
a theory is being tested when it is not. \emph{Innovation} is a property
of the product, not of the sentence describing it. What
\emph{atypicality} says is narrower than any of these: this description
is unusual for its class and year. Whether the thing described was new,
radical, or good is not measured.
